\section{Baselines}
\label{sec:app_baselines}
% =============================================================================

% =============================================================================
\subsection{Frontier LLMs with Verbalized Tools}
\label{sec:app_baseline_frontier}
% =============================================================================

To select the strongest frontier baseline, we evaluate five
models---GPT-5.1~GPT-5.1, Claude Sonnet~4.5,
Claude Opus~4.5, Gemini~2.5~Pro, and
Qwen3-235B~\cite{yang2025qwen3}---each given access to Lc0-BT4 via a
ReAct~\cite{yao2022react} tool-calling loop.
At each invocation the tool returns the top-5 moves with centipawn
evaluations, win/draw/loss probabilities, and principal variations up to
depth~20.
All models share identical tool schemas, system prompts, and sampling
parameters; the only variable is the LLM backbone.
We sample 100~positions from each LLAMIA-Bench task and report the
aggregate metric per task group.

%% ── TABLE: FRONTIER MODEL SELECTION ─────────────────────────────────────────
\begin{table*}[ht]
\centering
\setlength{\tabcolsep}{5pt}
\caption{%
  \textbf{Frontier model selection.}
  Average metric per LLAMIA-Bench task group, all models using Lc0-BT4
  verbalized tool access via ReAct.
  Per-group scores are unweighted means over the constituent columns of
  \cref{tab:full_bench}: BC = mean(MAIA, Wild); Puzzle = mean(Difficulty,
  Interest); Annot.\ = Rationale; Comm.\ = Commentary.
  The GPT-5.1\,+\,Lc0 row aggregates directly from \cref{tab:full_bench};
  the remaining rows are evaluated on a matched 100-sample subset per task
  with identical tool schemas, prompts, and sampling parameters.
  Gameplay is excluded because it requires the full gauntlet protocol.
  GPT-5.1 achieves the highest aggregate and is adopted as the frontier
  verbalized baseline in all subsequent experiments.
}
\label{tab:frontier_selection}
\begin{adjustbox}{max width=\textwidth}
\renewcommand{\arraystretch}{1.3}
\small
\begin{tabular}{lccccc}
\toprule[1.2pt]
\textbf{Model} & \textbf{BC\,$\uparrow$} & \textbf{Puzzle\,$\uparrow$}
  & \textbf{Annot.\,$\uparrow$} & \textbf{Comm.\,$\uparrow$}
  & \textbf{Avg.\,$\uparrow$} \\
\midrule
GPT-5.1\,+\,Lc0
  & \valbest{42.5} & \valbest{29.0} & \valbest{37.5}
  & \valbest{52.0} & \valbest{40.3} \\
Claude Opus~4.5\,+\,Lc0
  & 39.6 & 27.4 & 36.3 & 50.3 & 38.4 \\
Claude Sonnet~4.5\,+\,Lc0
  & 36.9 & 26.3 & 34.2 & 47.1 & 36.1 \\
Gemini~2.5~Pro\,+\,Lc0
  & 40.6 & 28.1 & 36.7 & 49.8 & 38.8 \\
Qwen3-235B\,+\,Lc0
  & 36.5 & 25.5 & 33.4 & 45.6 & 35.3 \\
\bottomrule[1.2pt]
\end{tabular}
\end{adjustbox}
\end{table*}

GPT-5.1 obtains the highest average across all task groups.
We therefore use \textbf{GPT-5.1\,+\,Lc0\,(Verb)} as the frontier
verbalized baseline throughout the paper.

% =============================================================================
\subsection{LLAMIA-Verb}
\label{sec:app_baseline_verb}
% =============================================================================

LLAMIA-Verb is the primary controlled ablation of LLAMIA.
It receives the identical base model, training data, reward signals, and RL
recipe (DAPO) as LLAMIA, but the subagent's output is provided exclusively
through verbalized tool responses: top-$k$ moves, centipawn evaluations, WDL
probabilities, and principal variations rendered as text tokens.
No LatentBridge projection is trained; the continuous latent
tokens~$z_S$ that LLAMIA receives are replaced by their textual equivalents.

We train LLAMIA-Verb at three scales---4B, 8B, and 14B---matching the
corresponding LLAMIA checkpoints in base model architecture, training data,
and total compute budget.
This controlled setup isolates the contribution of latent state
internalization from model family, data mix, reward shaping, and
optimization, and directly tests the central claim that verbalization is a
lossy bottleneck (\cref{sec:experiments}).

% =============================================================================
\subsection{LLAMIA (4B, 8B, 14B)}
\label{sec:app_baseline_llamia_scale}
% =============================================================================

To test whether the verbalization debt is an artifact of scale rather than interface, we train LLAMIA at three scales---4B, 8B, and 14B---matching the corresponding LLAMIA-Verb checkpoints in base model architecture, training data, and total compute budget.
All three LLAMIA checkpoints use the full latent-state internalization pipeline: LatentBridge projection of BT4 activations into $k{=}32$ continuous tokens, Stage~1 projector alignment, and Stage~2 end-to-end DAPO.
Comparing LLAMIA-$n$B against LLAMIA-Verb-$n$B at each scale isolates the interface contribution (internalization vs.\ verbalization) independently of model capacity, and directly supports the claim that the performance gap is not closed by scaling the LLM (\cref{sec:experiments}).

% =============================================================================
\subsection{Dedicated task finetunes}
\label{sec:app_baseline_experts}
% =============================================================================

For each LLAMIA-Bench task we compare against the strongest published or
reproducible task-specific model.
Table~\ref{tab:task_experts} lists the expert per task alongside its training
paradigm and data scale.
These models represent the performance ceiling achievable with task-specific
architectures and, in several cases, substantially more training data than
LLAMIA receives.
Tasks marked~--- have no established prior expert; LLAMIA-Bench introduces
them as new evaluation targets.

%% ── TABLE: TASK EXPERTS ─────────────────────────────────────────────────────
\begin{table*}[ht]
\centering
\setlength{\tabcolsep}{5pt}
\caption{%
  \textbf{Task experts used in LLAMIA-Bench evaluation.}
  Each row lists the strongest available dedicated task finetune for a given task.
  Tasks marked~--- are new evaluation targets with no prior
  task-specific model.
  \itwmark{} denotes in-the-wild splits.
}
\label{tab:task_experts}
\begin{adjustbox}{max width=\textwidth}
\renewcommand{\arraystretch}{1.35}
\small
\begin{tabular}{%
  p{3.8cm}
  p{4.0cm}
  p{1.8cm}
  p{2.0cm}
}
\toprule[1.2pt]
\textbf{Task / Split} & \textbf{Task Expert} & \textbf{Method}
  & \textbf{Data Scale} \\
\midrule


%% ── BEHAVIOR CLONING ────────────────────────────────────────────────────────
\rowcolor{clrBC}
\multicolumn{4}{l}{\enspace\small\textbf{Behavior Cloning}} \\
\rowcolor{clrBC}
Maia (Elo buckets)
  & Allie~\cite{zhang2024humanalignedchessbitsearch}
  & SL+Search
  & 93M games \\
\rowcolor{clrBC}
GM-25\,/\,Low-Time\,/\,Elo Gap\,\itwmark{}
  & Allie~\cite{zhang2024humanalignedchessbitsearch}
  & SL+Search
  & 93M games \\
\midrule


%% ── PUZZLE UNDERSTANDING ────────────────────────────────────────────────────
\rowcolor{clrPuzzle}
\multicolumn{4}{l}{\enspace\small\textbf{Puzzle Understanding}} \\
\rowcolor{clrPuzzle}
Difficulty Estimation
  & \cite{milosz2024predicting}
  & SFT
  & 4M \\
\rowcolor{clrPuzzle}
Interest Estimation
  & ---
  & ---
  & --- \\
\midrule

%% ── MOVE ANNOTATION ─────────────────────────────────────────────────────────
\rowcolor{clrAnnot}
\multicolumn{4}{l}{\enspace\small\textbf{Move Annotation}} \\
\rowcolor{clrAnnot}
Move Annotation
  & SCC~\cite{zang2019automated}
  & SL
  & 90K games \\

\bottomrule[1.2pt]
\end{tabular}
\end{adjustbox}
\end{table*}

% =============================================================================
